\documentclass[a4paper,11pt]{article}

\usepackage[utf8]{inputenc}
\usepackage[T1]{fontenc}
\usepackage[ngerman]{babel}
\usepackage{amsmath,amssymb}
\usepackage[most]{tcolorbox}
\usepackage{microtype}
\usepackage{geometry}
\usepackage{enumitem}
\usepackage{hyperref}

\geometry{margin=2.3cm}
\setlength{\parindent}{0pt}
\setlength{\parskip}{0.55em}
\setlist{nosep}

\newtcolorbox{mathbox}[1]{
	colback=gray!6,
	colframe=gray!55,
	title={#1},
	fonttitle=\bfseries,
	arc=1mm,
	boxrule=0.5pt,
	left=2mm,
	right=2mm,
	top=1mm,
	bottom=1mm
}

\title{\textbf{Moral und Theismus}\\
	\large Eine philosophische Untersuchung des Zusammenhangs
	zwischen objektiver Moral, Naturalismus und Gottesexistenz}
\author{}
\date{}

\begin{document}
	
	\maketitle
	
	\begin{abstract}
		Die Frage, ob moralische Urteile objektiv gültig sein können, bildet
		einen zentralen Ausgangspunkt für die Diskussion über das Verhältnis
		von Moral, Naturalismus und Theismus. Der hier untersuchte Ansatz
		argumentiert, dass ein materialistisch-naturalistisches Weltbild
		Schwierigkeiten hat, die Objektivität moralischer Tatsachen zu erklären.
		Wenn moralische Überzeugungen vorwiegend durch evolutionäre Prozesse
		entstanden sind, folgt daraus nicht unmittelbar, dass sie objektiv
		wahr sind. Wird zugleich vorausgesetzt, dass objektive moralische
		Tatsachen existieren, ergibt sich ein Argument gegen den evolutionären
		Naturalismus. Durch eine zusätzliche Annahme über das Verhältnis von
		Gottesexistenz und Naturalismus wird daraus ein Argument für den
		Theismus. Die Überzeugungskraft dieses Schlusses hängt jedoch von
		mehreren umstrittenen Voraussetzungen ab, insbesondere vom moralischen
		Realismus und von der Annahme, dass ohne Gott keine objektive Moral
		bestehen kann.
	\end{abstract}
	
	\section{Einleitung}
	
	Moralische Urteile prägen menschliches Denken und Handeln. Menschen
	bewerten Handlungen als richtig oder falsch, gut oder böse, unabhängig
	davon, ob diese Bewertungen für sie persönlich vorteilhaft sind.
	So erscheint beispielsweise die Folter unschuldiger Kinder nicht
	lediglich als unerwünscht, sondern als grundsätzlich falsch. Daraus
	ergibt sich eine philosophische Frage: Sind moralische Urteile nur
	Ausdruck individueller oder gesellschaftlicher Einstellungen, oder
	gibt es moralische Tatsachen, die unabhängig von solchen Einstellungen
	gelten?
	
	Diese Frage ist für die Weltanschauung von Bedeutung. Wenn moralische
	Werte objektiv existieren, stellt sich die weitere Frage, worauf ihre
	Gültigkeit beruht. Eine mögliche Antwort sieht ihr Fundament in einer
	göttlichen Wirklichkeit. Eine andere Möglichkeit besteht darin,
	Moral ohne Bezug auf eine übernatürliche Instanz zu erklären.
	
	Die untersuchte Argumentation verfolgt deshalb zwei Ziele. Erstens
	soll gezeigt werden, dass der evolutionäre Naturalismus Schwierigkeiten
	hat, die Existenz objektiver moralischer Tatsachen zu begründen.
	Zweitens soll aus dieser Kritik zusammen mit einer zusätzlichen
	Annahme ein Argument für die Existenz Gottes entwickelt werden.
	Dabei wird ein materialistisch-naturalistisches Weltbild betrachtet,
	in dem alle grundlegenden Vorgänge auf materielle Prozesse und
	Naturgesetze zurückgeführt werden.
	
	\section{Grundbegriffe}
	
	\subsection{Naturalismus und Evolution}
	
	Der materialistische Naturalismus geht davon aus, dass die Wirklichkeit
	grundsätzlich durch Materie und physikalische Gesetzmässigkeiten
	bestimmt ist. Nichtmaterielle Entitäten, die unabhängig von diesen
	Prozessen auf die physische Welt einwirken, sind in diesem Weltbild
	nicht vorgesehen.
	
	Der evolutionäre Naturalismus ergänzt diese Sichtweise durch eine
	evolutionäre Erklärung des Lebens. Die Entstehung und Entwicklung
	biologischer Organismen wird dabei insbesondere durch Mutation,
	Rekombination, genetische Drift und natürliche Selektion erklärt.
	Mutationen und andere genetische Veränderungen erzeugen Variation,
	während natürliche Selektion dazu führt, dass bestimmte Merkmale
	unter gegebenen Umweltbedingungen häufiger weitergegeben werden.
	
	Auch menschliche Denkweisen und Verhaltensdispositionen können im
	Rahmen dieser Perspektive als Ergebnisse evolutionärer Prozesse
	verstanden werden. Dazu könnten soziale Verhaltensweisen wie
	Kooperation, Fürsorge und die Ablehnung bestimmter Formen von Gewalt
	gehören. Ob daraus auch eine Erklärung der Gültigkeit moralischer
	Urteile folgt, ist allerdings eine eigenständige philosophische Frage.
	
	\subsection{Objektive Moral und Theismus}
	
	Objektive moralische Tatsachen sind Aussagen über richtiges oder
	falsches Handeln, deren Wahrheitswert nicht davon abhängt, ob
	einzelne Menschen oder Gesellschaften sie anerkennen. Ihre Gültigkeit
	wäre demnach nicht allein durch persönliche Vorlieben oder kulturelle
	Konventionen bestimmt.
	
	Demgegenüber hängt die Gültigkeit subjektiver moralischer Urteile
	von den Einstellungen oder Wertmassstäben bestimmter Personen
	oder Gruppen ab. Moralischer Realismus bezeichnet in diesem
	Zusammenhang die Position, dass zumindest einige moralische
	Aussagen objektiv wahr sein können. Moralischer Antirealismus und
	moralischer Relativismus bestreiten beziehungsweise relativieren
	eine solche objektive Gültigkeit.
	
	Der hier untersuchte Theismus versteht Gott als ein metaphysisch
	notwendiges, nichtmaterielles Wesen. Gott soll dabei nicht bloss
	als Urheber der Welt, sondern auch als mögliches Fundament
	objektiver moralischer Tatsachen verstanden werden. Die moralische
	Gültigkeit wäre dann in Gottes Wesen begründet und nicht allein
	in wechselnden menschlichen Überzeugungen.
	
	\section{Das moralische Argument gegen den Naturalismus}
	
	\subsection{Die Grundidee}
	
	Das erste Argument, das auf Überlegungen von Mark D. Linville
	aufbaut, richtet sich gegen den evolutionären Naturalismus. Sein
	Ausgangspunkt ist die Annahme, dass moralische Überzeugungen im
	Rahmen eines evolutionären Weltbildes zumindest teilweise durch
	natürliche Selektion entstanden sind.
	
	Viele moralische Verhaltensweisen lassen sich tatsächlich als
	potenziell vorteilhaft für das Überleben und die Fortpflanzung
	interpretieren. Fürsorge gegenüber Nachkommen kann beispielsweise
	deren Überlebenschancen erhöhen. Kooperation kann den Zusammenhalt
	einer Gruppe fördern, während eine Einschränkung innergruppaler
	Gewalt unter bestimmten Bedingungen die Stabilität sozialer
	Gemeinschaften verbessert.
	
	Aus solchen Überlegungen ergibt sich jedoch ein Problem: Ein
	evolutionärer Vorteil erklärt zunächst, weshalb sich eine bestimmte
	Verhaltensweise oder Überzeugung entwickelt und verbreitet hat.
	Er zeigt nicht ohne Weiteres, dass die entsprechende moralische
	Überzeugung objektiv wahr ist.
	
	Eine Überzeugung kann für das Überleben nützlich sein, ohne dass
	ihre Nützlichkeit bereits ihre Wahrheit begründet. Wenn moralische
	Überzeugungen vor allem deshalb entstanden sind, weil sie sich
	evolutionär bewährt haben, muss daher zusätzlich erklärt werden,
	weshalb ihre Inhalte mit objektiven moralischen Tatsachen
	übereinstimmen sollten.
	
	\begin{mathbox}{Argumentationsschema: Gegen den evolutionären Naturalismus}
		\begin{enumerate}[label=(\arabic*)]
			\item Wenn der evolutionäre Naturalismus wahr ist, sind moralische
			Überzeugungen Produkte natürlicher Selektion.
			\item Wenn moralische Überzeugungen Produkte natürlicher Selektion
			sind, existieren keine objektiven moralischen Tatsachen.
			\item Es existieren objektive moralische Tatsachen.
			\item Daher ist der evolutionäre Naturalismus falsch.
		\end{enumerate}
		
		\textit{Logische Struktur:} Aus den drei Prämissen folgt die
		Zurückweisung des evolutionären Naturalismus. Die formale
		Gültigkeit des Schlusses bedeutet allerdings nicht, dass seine
		Prämissen tatsächlich wahr sind.
	\end{mathbox}
	
	\subsection{Die Entstehung moralischer Überzeugungen}
	
	Für die erste Prämisse wird angenommen, dass menschliche geistige
	und körperliche Eigenschaften im materialistischen Naturalismus
	letztlich durch materielle Prozesse bestimmt werden. Wenn auch
	moralische Überzeugungen zu diesen Eigenschaften gehören, können
	sie als Ergebnisse biologischer und evolutionärer Entwicklungen
	verstanden werden.
	
	Dafür spricht, dass zahlreiche moralische Verhaltensweisen mit
	Überlebens- und Fortpflanzungsvorteilen verbunden sein können.
	Elterliche Fürsorge ist hierfür ein naheliegendes Beispiel.
	Auch Kooperation und die Einschränkung schädigenden Verhaltens
	können unter bestimmten Bedingungen die Stabilität einer Population
	begünstigen.
	
	Diese Beispiele machen eine evolutionäre Erklärung moralischer
	Dispositionen plausibel. Sie beweisen jedoch nicht, dass jede
	moralische Überzeugung evolutionär entstanden ist oder dass
	natürliche Selektion allein die Entstehung menschlicher Moral
	vollständig erklärt. Die Argumentation ist insofern als induktiver
	Schluss zu verstehen: Aus einzelnen plausiblen Beispielen wird
	auf einen allgemeineren Zusammenhang geschlossen.
	
	\begin{mathbox}{Mathematische Einordnung: Bayesianisches Argumentieren}
		Das Bayes'sche Theorem beschreibt, wie sich die Plausibilität
		einer Hypothese angesichts neuer Informationen verändert. In
		der sogenannten Odds-Form gilt:
		
		\[
		\frac{P(H\mid E)}{P(\neg H\mid E)}
		=
		\frac{P(H)}{P(\neg H)}
		\cdot
		\frac{P(E\mid H)}{P(E\mid\neg H)}.
		\]
		
		Dabei bezeichnet \(H\) eine Hypothese und \(E\) eine beobachtete
		Evidenz. Die erste Verhältnisgrösse beschreibt die anfängliche
		Plausibilität der Hypothese, die zweite den Einfluss der Evidenz.
		
		Auf die Evolution moralischer Überzeugungen übertragen bedeutet
		dies: Wenn moralische Dispositionen unter bestimmten Bedingungen
		einen Selektionsvorteil begünstigen, kann dies ihre evolutionäre
		Entstehung plausibler machen. Die Gleichung liefert jedoch nur
		dann eine belastbare Begründung, wenn die verwendeten
		Wahrscheinlichkeiten unabhängig und sachlich gerechtfertigt
		abgeschätzt werden können. Sie beweist für sich genommen weder
		die evolutionäre Entstehung aller Moralvorstellungen noch deren
		fehlende Objektivität.
	\end{mathbox}
	
	\subsection{Vom evolutionären Ursprung zur fehlenden Objektivität}
	
	Die zweite Prämisse ist philosophisch anspruchsvoller. Sie beruht
	auf der sogenannten Unabhängigkeitsthese: Die Entstehung einer
	Überzeugung und ihre Wahrheit sind zwei unterschiedliche Fragen.
	Aus der Tatsache, dass eine Überzeugung durch einen bestimmten
	Prozess entstanden ist, folgt nicht unmittelbar, ob sie wahr oder
	falsch ist.
	
	Ein einfaches Beispiel verdeutlicht diesen Punkt. Eine Person
	könnte durch eine zufällige Ziehung zu der Überzeugung gelangen,
	dass sich eine bestimmte Anzahl von Studierenden in einem
	Hörsaal befindet. Die Ziehung erklärt, wie die Person zu ihrer
	Überzeugung gelangt ist. Sie liefert aber keinen zuverlässigen
	Grund dafür, dass die angegebene Zahl tatsächlich stimmt.
	
	Ähnlich könnte die Evolution bestimmte moralische Überzeugungen
	begünstigt haben, ohne dass der evolutionäre Prozess auf die
	Wahrheit dieser Überzeugungen ausgerichtet war. Natürliche
	Selektion begünstigt Merkmale, die unter gegebenen Bedingungen
	den Fortpflanzungserfolg erhöhen. Daraus folgt nicht automatisch,
	dass sie Überzeugungen hervorbringt, die objektive moralische
	Tatsachen korrekt erfassen.
	
	Ein Gedankenexperiment soll diese Schwierigkeit verdeutlichen:
	Unter unterschiedlichen Umweltbedingungen könnten unterschiedliche
	Verhaltensweisen vorteilhaft sein. In einer Situation mit extremer
	Knappheit könnten Verhaltensdispositionen begünstigt werden, die
	in einer anderen Umweltgemeinschaft keinen Vorteil bieten oder
	sogar schädlich wären. Evolutionär entstandene Moralvorstellungen
	könnten daher von den jeweiligen Lebensbedingungen beeinflusst
	werden.
	
	Das Gedankenexperiment soll zeigen, dass ein evolutionärer
	Ursprung allein keine hinreichende Begründung für objektive
	moralische Gültigkeit darstellt. Es beweist jedoch nicht, dass
	evolutionär entstandene Überzeugungen grundsätzlich falsch sind.
	Ebenso wenig folgt aus unterschiedlichen evolutionären Ursprüngen
	unmittelbar, dass es keine objektiven moralischen Tatsachen gibt.
	Dafür wäre eine zusätzliche Begründung erforderlich.
	
	\subsection{Die Annahme objektiver Moral}
	
	Die dritte Prämisse setzt voraus, dass objektive moralische Tatsachen
	existieren. Als intuitiv überzeugendes Beispiel dient die Annahme,
	dass die Folter unschuldiger Kinder objektiv falsch ist und nicht
	bloss deshalb abgelehnt wird, weil eine bestimmte Gesellschaft
	sie missbilligt.
	
	Diese Überlegung ist jedoch keine abschliessende Begründung des
	moralischen Realismus. Dass eine moralische Aussage vielen
	Menschen selbstverständlich erscheint, beweist nicht, dass sie
	unabhängig von menschlichen Einstellungen wahr ist. Wer einen
	moralischen Antirealismus oder Relativismus vertritt, kann deshalb
	gerade diese Prämisse zurückweisen.
	
	Das Argument gegen den Naturalismus ist folglich nur dann
	überzeugend, wenn sowohl die evolutionäre Erklärung moralischer
	Überzeugungen als auch der behauptete Zusammenhang zwischen
	deren Entstehung und fehlender Objektivität akzeptiert werden.
	Der moralische Realismus bildet dabei eine grundlegende
	Voraussetzung.
	
	\section{Das moralische Argument für den Theismus}
	
	\subsection{Vom Naturalismus zur Gottesexistenz}
	
	Im zweiten Schritt wird das bisherige Argument um eine weitere
	Prämisse ergänzt. Diese besagt, dass der evolutionäre Naturalismus
	zutrifft, wenn kein göttliches Wesen existiert. Die Begründung
	dafür lautet, dass keine hinreichend plausible Alternative zur
	naturalistischen Erklärung der Entstehung und Entwicklung des
	Lebens angenommen wird.
	
	Wird diese Zusatzannahme akzeptiert, lässt sich die Argumentation
	folgendermassen erweitern: Wenn objektive moralische Tatsachen
	existieren, ist der evolutionäre Naturalismus nach dem vorherigen
	Argument falsch. Wenn zugleich gilt, dass ohne Gott der
	evolutionäre Naturalismus wahr wäre, kann aus der Falschheit
	des evolutionären Naturalismus auf die Existenz Gottes
	geschlossen werden.
	
	\begin{mathbox}{Argumentationsschema: Für den Theismus}
		\begin{enumerate}[label=(\arabic*)]
			\item Der evolutionäre Naturalismus ist mit der Existenz
			objektiver moralischer Tatsachen unvereinbar.
			\item Es existieren objektive moralische Tatsachen.
			\item Daher ist der evolutionäre Naturalismus falsch.
			\item Wenn kein Gott existiert, ist der evolutionäre
			Naturalismus wahr.
			\item Daher existiert Gott.
		\end{enumerate}
		
		Die Schlussfolgerung folgt logisch aus den Prämissen. Ihre
		Überzeugungskraft hängt aber insbesondere davon ab, ob die
		vierte Prämisse gerechtfertigt ist und ob die vorherigen
		Prämissen tatsächlich zutreffen.
	\end{mathbox}
	
	\subsection{Gott als Fundament der Moral}
	
	Die zugrunde liegende Vorstellung ist, dass objektive moralische
	Tatsachen in Gottes Wesen begründet sein könnten. Moralische
	Gültigkeit wäre dann nicht von wechselnden gesellschaftlichen
	Normen oder individuellen Vorlieben abhängig. Sie würde vielmehr
	auf einem metaphysisch notwendigen Fundament beruhen.
	
	Damit ist allerdings noch nicht gezeigt, dass die objektive Moral
	tatsächlich in Gott begründet ist. Selbst wenn der evolutionäre
	Naturalismus als unzureichende Erklärung moralischer Objektivität
	zurückgewiesen wird, folgt daraus allein noch nicht, dass jede
	andere naturalistische oder nichttheistische Position ebenfalls
	falsch ist. Der Schluss auf Gott setzt voraus, dass die
	vorgeschlagene Alternative hinreichend umfassend ist.
	
	\subsection{Das Euthyphron-Dilemma}
	
	Ein wichtiger Einwand gegen die Vorstellung, Gott begründe die
	Moral, ist das Euthyphron-Dilemma. Es fragt, ob etwas moralisch
	gut ist, weil Gott es will, oder ob Gott es will, weil es bereits
	moralisch gut ist.
	
	Im ersten Fall scheint Moral vom göttlichen Willen abzuhängen.
	Dadurch stellt sich die Frage, ob moralische Normen auch anders
	sein könnten, wenn Gott etwas anderes gewollt hätte. Im zweiten
	Fall scheint das moralisch Gute unabhängig von Gott zu bestehen.
	Dann wäre Gott nicht das eigentliche Fundament der moralischen
	Wahrheit.
	
	Als mögliche Antwort wird vorgeschlagen, das Gute nicht als
	willkürliches Ergebnis einzelner göttlicher Entscheidungen zu
	verstehen, sondern es in Gottes Wesen selbst zu verankern.
	Moralische Tatsachen wären demnach gut, weil sie mit dem Wesen
	Gottes übereinstimmen. Diese Lösung soll sowohl die Willkür des
	ersten Ansatzes als auch die Unabhängigkeit der Moral vom
	zweiten Ansatz vermeiden. Ob sie das Dilemma vollständig löst,
	bleibt jedoch eine eigenständige philosophische Streitfrage.
	
	\section{Schlussfolgerung}
	
	Die untersuchte Argumentation verbindet drei zentrale Gedanken:
	Erstens können moralische Überzeugungen als Ergebnisse
	evolutionärer Prozesse verstanden werden. Zweitens erklärt
	eine solche Entstehungsgeschichte noch nicht, weshalb diese
	Überzeugungen objektiv wahr sein sollten. Drittens könnte die
	Annahme objektiver moralischer Tatsachen deshalb als Einwand
	gegen einen materialistisch verstandenen evolutionären
	Naturalismus dienen.
	
	Wird zusätzlich angenommen, dass ohne Gott der evolutionäre
	Naturalismus zutrifft, ergibt sich daraus ein Argument für den
	Theismus. Die Schlussfolgerung ist jedoch nicht unabhängig von
	ihren Voraussetzungen. Insbesondere die Existenz objektiver
	moralischer Tatsachen, die fehlende Verbindung zwischen
	evolutionärem Ursprung und moralischer Wahrheit sowie die
	Annahme, dass ohne Gott keine plausible Alternative zum
	Naturalismus besteht, sind philosophisch umstritten.
	
	Entscheidend ist daher die Unterscheidung zwischen einem logisch
	gültigen Schluss und einem inhaltlich überzeugenden Argument.
	Die formale Gültigkeit zeigt, dass die Schlussfolgerung aus den
	Prämissen folgt. Sie zeigt nicht, dass diese Prämissen wahr
	sind. Das moralische Argument kann somit als ein Versuch
	verstanden werden, die Plausibilität des Theismus zu erhöhen,
	nicht jedoch als ein voraussetzungsloser Beweis der
	Gottesexistenz.
	
	\begin{thebibliography}{9}
		
		\bibitem{Craig}
		Craig, W. L. (2009).
		\textit{Reasonable Faith: Christian Truth and Apologetics}.
		3. Auflage. Crossway Books, Wheaton.
		
		\bibitem{Joyce}
		Joyce, R. (2006).
		\textit{The Evolution of Morality}.
		MIT Press, Cambridge, MA.
		
		\bibitem{Linville}
		Linville, M. D. (2009).
		The Moral Argument.
		In: W. L. Craig und J. P. Moreland (Hrsg.),
		\textit{The Blackwell Companion to Natural Theology},
		S. 391--448. Blackwell.
		
		\bibitem{Mawson}
		Mawson, T. (2009).
		Morality and Religion.
		\textit{Philosophy Compass}, 4(6), S. 1033--1043.
		
		\bibitem{Sober}
		Sober, E. (1994).
		\textit{From a Biological Point of View: Essays in Evolutionary Philosophy}.
		Cambridge University Press, Cambridge.
		
	\end{thebibliography}
	
\end{document}
```